\section{Related Work}

\paragraph{Efficient Tool Calling}
As large language models (LLMs) increasingly rely on multi-step, multi-turn tool-calling workflows, inference latency in LLM serving has grown substantially. To address this, recent studies leverage directed acyclic graphs (DAGs) to model inter-tool dependencies, enabling parallel execution of independent API calls and thereby reducing latency~\cite{LLMCompiler, zhu2025:dividethenaggregate}. Another line of work executes tool calls proactively during generation to minimize system idle time~\citep{xu2024conveyor, nichols2025spec}. In contrast to prior methods that focus on tool execution, this work shifts attention to accelerating tool-calling generation, further improving the inference efficiency of tool-use LLMs.

\paragraph{Speculative Decoding}
In recent years, speculative decoding (SD) has emerged as a widely adopted paradigm for accelerating LLM inference~\cite{leviathan:2023spec, xia:2023specdec}. At each decoding step, SD first generates a sequence of draft tokens, which are subsequently verified in parallel by the target LLM~\citep{xia:2024survey}. Recent advances in SD, such as the Eagle series~\citep{eagle, li2024:eagle2, li2025:eagle3}, introduce lightweight draft models layered on top of the target LLM, achieving promising speedups. While these approaches require additional training or parameters, an alternative line of research focuses on training-free SD methods~\citep{luo:2025tokenrecycle, saxena2023pld, SWIFT}. Notably, Token Recycling~\citep{luo:2025tokenrecycle} stores the top-$k$ candidate tokens of previous steps in an adjacency matrix and uses breadth-first search to construct a draft tree. Similarly, SAM-Decoding~\citep{hu2025:samd} leverages a suffix automaton to efficiently retrieve high-quality drafts from both the input context and a static text corpus.